Ce projet a pour objectif de concevoir un robot autonome capable de sortir d’un labyrinthe. 
Il s’inscrit dans une démarche visant à approfondir nos connaissances en algorithmique, en 
programmation en langage C, et en électronique, tout en nous initiant aux pratiques collaboratives, 
essentielles dans le cadre de notre formation d’ingénieur.

Le robot sera doté de capteurs variés lui permettant de collecter des informations sur son
environnement et d’exécuter les instructions programmées à l’avance. À partir d’un plan de 
labyrinthe donné en entrée, les algorithmes développés auront pour mission de générer un chemin 
vers la sortie, ainsi que les ordres nécessaires à l’exécution de ce trajet par le 
robot.

Ce compte rendu présente les différentes étapes du projet : une première partie sera consacrée à 
la conception algorithmique et électronique, tandis qu’une seconde partie abordera les 
problématiques rencontrées au cours de ces phases.